% ML-only related work — citations from filtered ML reference set
\section{Related Work}
\label{sec:related}

\subsection{Handcrafted texture features for MRI}
Grey-level co-occurrence matrices (GLCM), local binary patterns (LBP), histogram of oriented gradients (HOG), and discrete wavelet transforms (DWT) encode complementary tumour micro-texture~\cite{ref_fonc20231248452,ref_hsr271195,ref_s44163024002144}.
Yu \textit{et al.}~\cite{ref_fnins2021634926} showed that texture features with an SVM classifier differentiate tumour from healthy regions on MRI with high sensitivity.
Comparative studies consistently rank texture and gradient descriptors among the strongest classical representations~\cite{ref_ijigsp20150303,ref_iasc2023032426}.
Our Phase~1 differs by selecting parameters through \emph{classifier-free} separability metrics before any model is trained.

\subsection{Classical ML classifiers in brain tumour classification}
Support vector machines with RBF kernels remain a dominant baseline on radiomic feature vectors~\cite{ref_s0033002007108w,ref_s41598024772437}.
Ensemble tree methods provide native feature-importance scores for interpretability~\cite{ref_iasc2023032426,ref_jdaip202084020}.
Prior benchmarks rarely compare more than three classifiers on identical handcrafted features and splits.

\subsection{Parameter selection strategies}
Default scikit-learn parameters are common; grid search coupled to a single classifier is expensive and risks overfitting the validation objective to one model family.
Separability metrics offer a lightweight alternative used in feature selection but seldom applied systematically across texture families for brain MRI.

\subsection{Gap}
No study combines classifier-free separability scoring for GLCM, LBP, DWT, and HOG with an eight-model benchmark and full statistical validation on the same dataset and preprocessing chain.
We fill this gap with a reproducible two-phase notebook pipeline on a fixed stratified 80/20 split.
